% section_perm.tex — Section 9: permutation monomials (first solver, 2026-08-19).  Self-contained; uses \Fp, \Z, amsthm environments,
% Theorem~\ref{thm:cubic}, Lemma~\ref{lem:projection}, Corollary~\ref{cor:cubic} of Section~\ref{sec:cubic}, \cite{IK,Perelmuter,Montgomery,Vaaler}.
% New references needed in the bibliography:
%   \bibitem{Wan93} D.~Wan, A $p$-adic lifting lemma and its applications to permutation polynomials, in: Finite Fields, Coding Theory, and Advances in
%     Communications and Computing (Las Vegas 1991), Lecture Notes in Pure and Appl.\ Math.\ 141, Dekker, 1993, 209--216.  [only for Remark]
%   \bibitem{NR82} H.~Niederreiter and K.~H. Robinson, Complete mappings of finite fields, J.\ Austral.\ Math.\ Soc.\ Ser.\ A 33 (1982) 197--212.  [only for Remark]

\section{Permutation monomials: the same cover, the same conclusion}\label{sec:perm}

Section~\ref{sec:cubic} used only the lines of slope $\pm1$ and the exact geometry of their point-profiles.  Nothing in that argument is
special to degree three except the final bookkeeping, and here we carry it out for the permutation monomials $f(x)=x^k$, $k$ odd, $\gcd(k,p-1)=1$
(for even $k$ there are none).  The result is the same as for cubics: the explicit $\pm1$-line cover has cost about $1.40\,N$ for every $k$, so no
permutation monomial graph reaches the HJSW value.  The proof isolates three ``local law'' statements (root counts, root positions, and the two
slopes) which we prove for $x^k$, and a general lift-accounting lemma; the constant is an explicit rational number.

Throughout, $k\ge3$ is odd and fixed, $p\nmid k(k-1)$, $\gcd(k,p-1)=1$, $B=[x_0,x_0+2p)\times[y_0,y_0+2p)$, $P=P(x^k,B)$ ($4p$ points, two per row and
per column), and for $\sigma\in\{+1,-1\}$ we write $g_\sigma(t)=t^k-\sigma t$, so that $Y-X\equiv g_+(x)$ and $X+Y\equiv g_-(x)\pmod p$ for the lifts of the
residue $x$.  For $c\in\Fp$ let $R_\sigma(c)=\{t: g_\sigma(t)=c\}$; every residue $x$ lies in exactly one $R_\sigma(c)$ for each $\sigma$, and we write
$k_\sigma(x)=|R_\sigma(g_\sigma(x))|$.  Classes and thresholds are exactly as in Section~\ref{sec:cubic}: with $u_t=(r_t-x_0)/p$ the position of $t$ in the
$x$-window, $\theta_+(c)=1-((c-y_0+x_0)\bmod p)/p$ and $\theta_-(c)=((c-x_0-y_0)\bmod p)/p$, a root $t\in R_+(c)$ has class $0$ iff $u_t\ge\theta_+(c)$ and a root
$t\in R_-(c)$ has class $0$ iff $u_t\le\theta_-(c)$; for a single residue $x$ the two thresholds are tied by the identity
\begin{equation}\label{eq:tie}
 \theta_-(g_-(x))\ \equiv\ 1-\theta_+(g_+(x))+2u_x\pmod 1 ,
\end{equation}
because $g_-(x)-g_+(x)=2x$.

\subsection*{The lift-accounting lemma}

\begin{lemma}[profiles and lifts]\label{lem:lifts}
Let $c$ be a residue of slope $\sigma$ with $|R_\sigma(c)|=k'$ roots, $n_0$ of class $0$ and $n_1=k'-n_0$ of class $1$.  The four consecutive lines of slope
$\sigma$ carrying the $4k'$ lifts of these roots have $n_0,\ k'+n_0,\ k'+n_1,\ n_1$ points.  For a residue $x$ with $k_\pm=k_\pm(x)$ and $n_\pm=n_\pm(x)$
the number of roots of its residue of slope $\pm$ lying in the class of $x$ (so $1\le n_\pm\le k_\pm$), the lines through the four lifts
$(\delta,\varepsilon)$ of $x$ have the following numbers of points ($\{a,b\}$ means the two values in some order for the two lifts of the pair):
\[
\begin{array}{c|c|c}
\text{lifts} & \text{slope }+1 & \text{slope }-1\\ \hline
(0,0),(1,1) & k_++n_+,\ k_++n_+ & \{n_-,\ 2k_--n_-\}\\
(0,1),(1,0) & \{2k_+-n_+,\ n_+\} & k_-+n_-,\ k_-+n_-
\end{array}
\]
Consequently, for the cover $W_4$ (weight $1$ on every $\pm1$-line with $\ge4$ points, weight $1$ on every point lying on none of them),
\begin{align}
 U(x)&=[k_++n_+\le3]\bigl([n_-\le3]+[2k_--n_-\le3]\bigr)\notag\\
 &\qquad+[k_-+n_-\le3]\bigl([n_+\le3]+[2k_+-n_+\le3]\bigr),\label{eq:Ux}\\
 L&=\sum_\sigma\sum_c\Bigl([n_0\ge4]+[k'+n_0\ge4]+[k'+n_1\ge4]+[n_1\ge4]\Bigr),\label{eq:Lperm}\\
 \alpha(x^k,B)&\ \le\ \mathrm{cost}(W_4)=2L+\sum_{x\in\Fp}U(x).\label{eq:costW4}
\end{align}
\end{lemma}
\begin{proof}
For slope $+1$ the lifts of a class-$0$ root $t$ have $Y-X=c'$ (twice, $\delta=\varepsilon$), $c'+p$ ($(0,1)$) and $c'-p$ ($(1,0)$), those of a class-$1$ root
$c'+p$ (twice), $c'+2p$ ($(0,1)$) and $c'$ ($(1,0)$), exactly as in Section~\ref{sec:cubic}; counting gives the profile $(n_0,k'+n_0,k'+n_1,n_1)$ on the lines
$c'-p,c',c'+p,c'+2p$, and reading off the line of each lift of $x$ gives the first column of the table (a class-$0$ root has its equal-shift lifts on the
line with $k'+n_0=k'+n_+$ points, its $(0,1)$-lift on the line with $k'+n_1=2k'-n_+$ points and its $(1,0)$-lift on the line with $n_0=n_+$ points; for a
class-$1$ root the same three numbers occur with the two mixed lifts exchanged).  For slope $-1$ the roles of ``equal'' and ``mixed'' lifts are exchanged
(the profile sits on $X+Y\in\{c'',c''+p,c''+2p,c''+3p\}$ with the two mixed lifts of a root on the central lines).  A lift is a singleton of $W_4$ iff both its
lines have $\le3$ points; summing over the two lifts of each kind gives \eqref{eq:Ux}, and \eqref{eq:costW4} is the usual duality (a lawful set meets a line
in $\le2$ points, a point in $\le1$).
\end{proof}
For $k'=3$ this is the bookkeeping of Section~\ref{sec:cubic} ($k_++n_+\le3$ iff $k_+=1$ or $k_+=2$ with a split pair; the bracket $[n_-\le3]+[2k_--n_-\le3]$
equals $2$ for a same triple and $1$ for a split one).  Note that for $k'\ge4$ both central lines are always rich ($k'+n_0,k'+n_1\ge4$).

\subsection*{The three local laws for $x^k$}

\begin{lemma}[root counts]\label{lem:rootcounts}
Let $p\nmid k(k-1)$.  For each $\sigma$ and $0\le j\le k$,
$\#\{c\in\Fp:\ |R_\sigma(c)|=j\}=P_j\,p+O_k(\sqrt p)$ with $P_j=\frac1{j!}\sum_{i=0}^{k-j}\frac{(-1)^i}{i!}$ (the proportion of permutations of $k$ letters with
$j$ fixed points; $P_{k-1}=0$).  In particular the number of residues $x$ with $k_\sigma(x)=j$ is $jP_j\,p+O_k(\sqrt p)$.
\end{lemma}
\begin{proof}
The map $t\mapsto g_\sigma(t)$ has the $k-1$ simple critical points $t^{k-1}=\sigma/k$ with the pairwise distinct critical values $t\,(\sigma/k-\sigma)$,
so its geometric monodromy group is generated by transpositions and is transitive ($g_\sigma(t)-c$ is irreducible over $\overline\Fp(c)$, being linear in
$c$): it is $S_k$, hence so is the arithmetic monodromy group.  Chebotarev's density theorem for the function field $\Fp(c)$ (with the Weil error term, the
covering curve having genus bounded in terms of $k$) gives the count of $c$ whose Frobenius has $j$ fixed points, i.e.\ $j$ rational roots.
\end{proof}

\begin{lemma}[root positions]\label{lem:positions}
Let $1\le j\le k$ and $\sigma\in\{\pm1\}$.  Over the ordered $j$-tuples $(t_1,\dots,t_j)$ of distinct roots of a common residue $c$ of slope $\sigma$ (one for each $c$
and each ordered choice of $j$ distinct elements of $R_\sigma(c)$), the vector $(u_{t_1},\dots,u_{t_j},\theta_\sigma(c))\in(\mathbb R/\mathbb Z)^{j+1}$ is equidistributed
with error $O_k(\sqrt p\log^{j+1}p)$ for box conditions, except that for $j=k$ the positions satisfy the one linear relation $\sum_iu_{t_i}\equiv-kx_0/p$
(and are equidistributed on that hyperplane).  Consequently, for a residue with $j<k$ roots the number $n_0$ of class-$0$ roots is $\mathrm{Bin}(j,\theta)$
with $\theta$ uniform, i.e.\ uniform on $\{0,\dots,j\}$, and for a random member $x$ of such a residue $n_\sigma(x)=1+\mathrm{Bin}(j-1,q_\sigma)$, where
$q_\sigma$ is the probability that a partner has the class of $x$: $q_+=1-\theta_+$ if $u_x\ge\theta_+$ and $q_+=\theta_+$ otherwise;
$q_-=\theta_-$ if $u_x\le\theta_-$ and $q_-=1-\theta_-$ otherwise.
\end{lemma}
\begin{proof}
The variety $V_j$ of ordered $j$-tuples of distinct roots of $g_\sigma(t)=c$ is a curve whose components correspond to the orbits of the geometric monodromy
group $S_k$ on such tuples; $S_k$ is $j$-transitive, so $V_j$ is absolutely irreducible.  For $(h_1,\dots,h_j,h_0)\ne0$ the function
$\sum_ih_it_i+h_0g_\sigma(t_1)$ is non-constant on $V_j$ when $j<k$: the fibre of $V_j\to V_{j-1}$ (forget $t_j$) has $k-j+1\ge2$ points, so $t_j$ is not a
rational function of $(t_1,\dots,t_{j-1})$; inductively all $h_i=0$ for $i\ge2$, and then $h_1t_1+h_0g_\sigma(t_1)$ is non-constant unless $h_1=h_0=0$.  For
$j=k$ the same argument leaves exactly the relation $\sum_it_i=0$ (the coefficient of $t^{k-1}$ in $g_\sigma$).  Bombieri's bound on the curve $V_j$
\cite{IK} gives $O_k(\sqrt p)$ for each such exponential sum, and the Fourier expansion of the $j+1$ interval indicators (or Selberg's polynomials
\cite{Montgomery,Vaaler}) gives the box-count with error $O_k(\sqrt p\log^{j+1}p)$.  The consequences follow by integrating in $\theta$:
$\int_0^1\binom j m\theta^m(1-\theta)^{j-m}\,d\theta=1/(j+1)$.
\end{proof}

\begin{lemma}[the two slopes]\label{lem:twoslopes}
Let $D_\pm(x)=k^k(x^k\mp x)^{k-1}\mp(k-1)^{k-1}$ (up to a constant, the discriminant of $g_\pm(t)-g_\pm(x)$ as a polynomial in $t$), and let $p$ not
divide $k(k-1)\cdot\mathrm{Res}(D_+,D_-)$.  Then, over $x\in\Fp$, the pairs $(k_+(x),k_-(x))$ are asymptotically independent with the laws of
Lemma~\ref{lem:rootcounts} (size-biased), and, given $(u_x,\theta_+(g_+(x)))$, the positions of the partners of $x$ on the two slopes are independent and
uniform, with the tie \eqref{eq:tie} for the second threshold; all with the error terms of Lemma~\ref{lem:positions}.
\end{lemma}
\begin{proof}
The Galois closures $L_\pm$ over $\Fp(x)$ of $g_\pm(t)-g_\pm(x)$ are ramified (over the affine $x$-line) exactly at the zeros of $D_\pm$; since
$\mathrm{Res}(D_+,D_-)\not\equiv0$, these finite branch loci are disjoint, so $L_+\cap L_-$ is unramified over the affine line and tame at $\infty$
(its degree divides $k!<p$), hence trivial: $L_+$ and $L_-$ are linearly disjoint and the fibre product $W$ of the two root varieties over the $x$-line is
absolutely irreducible.  Chebotarev on the compositum gives the joint law of the two Frobenius classes (independence of the counts), and Bombieri on $W$
gives the joint equidistribution of $x$, $g_+(x)$ and the partner positions of both slopes (no linear relation between them can hold, as in
Lemma~\ref{lem:positions}, by linear disjointness); the tie \eqref{eq:tie} is an identity, not a hypothesis.  For $k=3,5,7$ the resultant is a non-zero
integer with only finitely many prime factors (computed symbolically: for $k=5$ its prime factors below $5000$ are $2,5,7,97,127$).
\end{proof}

\subsection*{The theorem}

For the rencontres numbers $P_j$ of Lemma~\ref{lem:rootcounts} define
\begin{equation}\label{eq:Ck}
\begin{aligned}
 C_k\ :=\ &2\sum_\sigma\sum_{j}P_j\,\mathbb E_{n_0\sim\mathrm U\{0..j\}}\bigl[[n_0\ge4]+[j+n_0\ge4]+[2j-n_0\ge4]+[j-n_0\ge4]\bigr]\\
 &+\int_0^1\!\!\int_0^1\mathbb E\bigl[U(x)\mid u_x=u,\ \theta_+=1-g\bigr]\,du\,dg ,
\end{aligned}
\end{equation}
where inside the expectation $k_\pm$ have the size-biased laws $jP_j$, $n_\pm=1+\mathrm{Bin}(k_\pm-1,q_\pm)$ with $q_\pm$ as in
Lemma~\ref{lem:positions} and $\theta_-=1-\theta_++2u\bmod1$, and $U(x)$ is \eqref{eq:Ux}.  The integrand is a piecewise polynomial in $(u,g)$ (four
regions cut out by $g+2u=1$, $g+2u=2$, $u+g=1$), so $C_k$ is an explicit rational number:
\[
 C_3=\tfrac{11}4=2.75,\qquad C_5=\tfrac{28183}{10080}=2.79593\ldots,\qquad C_7=\tfrac{15265237}{5443200}=2.80446\ldots
\]
(\texttt{slack/pp\_constant\_exact.py}; for $k=3$ the two summands are $\tfrac12\cdot2$ and $\tfrac74$, as in Section~\ref{sec:cubic}).

\begin{theorem}[permutation monomials]\label{thm:perm}
Let $k\in\{3,5,7\}$, and let $p$ be a prime with $\gcd(k,p-1)=1$ and $p\nmid k(k-1)\mathrm{Res}(D_+,D_-)$.  Then for every $2p\times2p$ box $B$,
\[
 \alpha(x^k,B)\ \le\ \bigl(C_k+\varepsilon_k+o(1)\bigr)\,p,\qquad 0\le\varepsilon_k\le\frac2{(k-1)!}+\frac8{k!},\quad \varepsilon_3=0,
\]
as $p\to\infty$; numerically $C_5+\varepsilon_5\le2.946$ and $C_7+\varepsilon_7\le2.809$.  In particular $\alpha(x^k,B)\le(1.474+o(1))\cdot2p<\tfrac32\cdot2p$:
no permutation monomial graph of degree $3$, $5$ or $7$ reaches the HJSW value.  The same holds for every odd $k$ for which
$\mathrm{Res}(D_+,D_-)\ne0$ (we have verified this for $k\le7$ only).
\end{theorem}
\begin{proof}
By Lemma~\ref{lem:lifts}, $\alpha\le2L+\sum_xU(x)$.  The count $L$ is a sum over residues of a function of $(j,n_0)$; by Lemmas~\ref{lem:rootcounts}
and~\ref{lem:positions} the residues with $j<k$ roots contribute the first term of \eqref{eq:Ck} times $p$, up to $O_k(\sqrt p\log^{k+1}p)$; the residues
with $j=k$ roots have density $P_k=1/k!$ per slope, and for them only the two outer indicators $[n_0\ge4]+[n_1\ge4]$ can deviate from the model
(the central lines are rich for $k\ge4$), so $L$ deviates by at most $2\cdot2P_kp=4p/k!$ and $2L$ by at most $8p/k!$ — this is the second part of
$\varepsilon_k$.  For $\sum_xU(x)$: by Lemma~\ref{lem:twoslopes} the joint law of $(k_+,n_+,k_-,n_-)$ for a random residue $x$ is, up to the error terms, the
one described after \eqref{eq:Ck}, except for the residues $x$ lying in a $k$-root residue of some slope (density $kP_k=1/(k-1)!$ per slope), where the
partner positions obey the linear relation of Lemma~\ref{lem:positions}; for such $x$ the only factor of \eqref{eq:Ux} that can be affected is a single
$[n_\sigma\le3]$ (the other factor $[k_\sigma+n_\sigma\le3]$ vanishes for $k_\sigma=k\ge4$), so replacing it by $1$ changes $\sum_xU(x)$ by at most
$2/(k-1)!\cdot p$ — the first part of $\varepsilon_k$.  All remaining terms are the integral in \eqref{eq:Ck} times $p$, with error $O_k(\sqrt p\log^{k+1}p)$.
For $k=3$, $k$-root residues are exactly the three-root residues, and Section~\ref{sec:cubic} shows that only pairwise data enter, whence $\varepsilon_3=0$
and $C_3=11/4$.  The numerical values of $C_k$ follow from the exact integration of the piecewise polynomial integrand.
\end{proof}

\begin{remark}[numerics; the true constant; other bijections]
(a) The bookkeeping of Lemma~\ref{lem:lifts} evaluates $\mathrm{cost}(W_4)$ in time $O(p)$ without enumerating lines
(\texttt{slack/pp\_exact\_UL\_np.py}).  At $p\approx10^7$ (three primes, box $(0,0)$) it gives $\sum_xU(x)/p=1.7500\pm0.0002$ for $k=3$ (the exact $7/4$) and
$1.7631\pm0.0002$ for $k=5$, against the value $17767/10080=1.76260$ of the tied-threshold law in \eqref{eq:Ck} — the small excess is the $k$-root
correction $\varepsilon_5$, so the true asymptotic constant for $x^5$ is $2.797p\approx1.3985\,N$; a model with independent slopes would give
$1519/864=1.7581$, which the data exclude.  At $p\le887$ the cost of $W_4$ fluctuates by a few percent (for $x^5$: mean $1.388\,N$, range $1.36$--$1.41$
over seven primes; for $x^7$: mean $1.415\,N$; \texttt{docs/research/g\_agents/x5\_x7\_numerics.md}); the LP over all $\pm1$-lines is $\approx1.32$--$1.34\,N$
and the LP over all lines with $\ge3$ points $\approx1.07$--$1.10\,N$, as for cubics.
(b) Why the hyperbola $x\mapsto1/x$ is exactly $3/2$: its two slopes are maximally dependent — the partner of $x$ is $\mp1/x$ on the two slopes, the two
root pairs are the same conic up to sign, and the pair of slope $+1$ shares its centre iff the pair of slope $-1$ does not — so the independence of
Lemma~\ref{lem:twoslopes} fails completely, and no cover of this type can go below $3(p-1)$ (Theorem~\ref{thm:main}).  For $x^k$ the partners of the two
slopes are roots of two different curves with disjoint branch loci, which is exactly Lemma~\ref{lem:twoslopes}.
(c) For a general permutation polynomial $f$ of degree $k$ the same proof gives $\alpha\le(C(P^+,P^-)+\varepsilon+o(1))p$ with the root-count laws
$P^\pm$ of $f\mp x$ in place of the rencontres numbers, as soon as the two covers $t\mapsto f(t)\mp t$ have monodromy $S_k$ and disjoint finite branch
loci; and if $f-x$ is not a permutation polynomial, at least a positive proportion of its residues carry $\ge2$ roots (Wan's bound $|V(f-x)|\le p-(p-1)/k$
\cite{Wan93}), while for $p$ large in terms of $k$ no polynomial of degree $\ge2$ is a complete mapping \cite{NR82} — so a linear saving from
$\pm1$-lines is available for every permutation polynomial of bounded degree, though with a constant that depends on the family.
\end{remark}
